\documentclass[11pt]{article}

\usepackage[margin=1in]{geometry}
\usepackage[T1]{fontenc}
\usepackage{lmodern}
\usepackage{amsmath}
\usepackage{amssymb}
\usepackage{tikz-cd}
\usepackage{cleveref}

\Crefname{figure}{Figure}{Figures}
\crefname{figure}{Figure}{Figures}

\title{Commutative Diagrams in Homological Algebra:\\ Short Exact Sequences and Pullback Squares}
\author{A.\ Khoo}
\date{\today}

\begin{document}
\maketitle

\begin{abstract}
We illustrate two elementary constructions of homological algebra using
commutative diagrams rendered with the \texttt{tikz-cd} package. The first is the
defining diagram of a short exact sequence of abelian groups; the second is a
pullback (fibre product) square together with the triangle that witnesses its
universal property. The diagrams are drawn entirely by \TeX{} and require no
external graphics.
\end{abstract}

\section{Short exact sequences}

Let $\mathcal{A}$ be an abelian category. A sequence of objects and morphisms
\[
  0 \longrightarrow A \xrightarrow{\,f\,} B \xrightarrow{\,g\,} C \longrightarrow 0
\]
is a \emph{short exact sequence} when $f$ is a monomorphism, $g$ is an
epimorphism, and $\operatorname{im} f = \ker g$. Exactness at $B$ is the single
substantive condition; exactness at $A$ and at $C$ records, respectively, that
$f$ is injective and that $g$ is surjective. We display the same data as a
commutative diagram in \cref{fig:ses}, where the hooked arrow
$A \hookrightarrow B$ emphasises that $f$ is a monomorphism and the two-headed
arrow $B \twoheadrightarrow C$ that $g$ is an epimorphism.

\begin{figure}[ht]
  \centering
  \begin{tikzcd}[column sep=large, row sep=large]
    0 \arrow[r]
      & A \arrow[r, "f", hook]
      & B \arrow[r, "g", two heads] \arrow[d, "\pi"']
      & C \arrow[r]
      & 0 \\
    & & B/\operatorname{im} f \arrow[ur, "\bar{g}"', dashed]
  \end{tikzcd}
  \caption{A short exact sequence $0 \to A \to B \to C \to 0$. The induced map
  $\bar{g}\colon B/\operatorname{im} f \to C$ (dashed) is the isomorphism
  supplied by the first isomorphism theorem, factoring $g$ through the quotient
  projection $\pi$.}
  \label{fig:ses}
\end{figure}

The diagram in \cref{fig:ses} encodes the first isomorphism theorem: because
$\operatorname{im} f = \ker g$, the epimorphism $g$ descends along the quotient
projection $\pi\colon B \to B/\operatorname{im} f$ to an isomorphism
$\bar{g}$ with $g = \bar{g} \circ \pi$. Commutativity of the lower triangle is
exactly the relation $g = \bar{g}\,\pi$, and the dashed arrow signals that
$\bar{g}$ is the unique map making that triangle commute.

\section{Pullback squares}

Given morphisms $f\colon X \to Z$ and $g\colon Y \to Z$ with common codomain,
their \emph{pullback} (or fibre product) is an object $P = X \times_Z Y$ equipped
with projections $p\colon P \to X$ and $q\colon P \to Y$ satisfying
$f \circ p = g \circ q$, and universal among all such cones. Set-theoretically,
\[
  X \times_Z Y \;=\; \bigl\{\, (x,y) \in X \times Y : f(x) = g(y) \,\bigr\},
\]
with $p$ and $q$ the restrictions of the two coordinate projections. The
commutative square is shown in \cref{fig:pullback}; the small right-angle mark in
its upper-left corner is the conventional \texttt{tikz-cd} annotation for a
pullback.

\begin{figure}[ht]
  \centering
  \begin{tikzcd}[column sep=large, row sep=large]
    P \arrow[r, "p"] \arrow[d, "q"'] \arrow[dr, phantom, "\lrcorner", very near start]
      & X \arrow[d, "f"] \\
    Y \arrow[r, "g"']
      & Z
  \end{tikzcd}
  \caption{The pullback (fibre product) square $P = X \times_Z Y$, with
  $f \circ p = g \circ q$. The corner symbol marks the square as cartesian.}
  \label{fig:pullback}
\end{figure}

The universal property promised by \cref{fig:pullback} is a statement about a
triangle. Suppose $W$ is any object carrying maps $a\colon W \to X$ and
$b\colon W \to Y$ with $f \circ a = g \circ b$. Then there is a unique mediating
morphism $u\colon W \to P$ through which both $a$ and $b$ factor, so that
$p \circ u = a$ and $q \circ u = b$. We isolate the mediating triangle for the
$X$-component in \cref{fig:universal}, where the natural transformation
$\alpha \Rightarrow \beta$ between the two composite functors
$W \rightrightarrows X$ is drawn as a two-cell.

\begin{figure}[ht]
  \centering
  \begin{tikzcd}[column sep=huge, row sep=large]
    W \arrow[r, "a"] \arrow[dr, "u"', dashed]
      \arrow[r, ""{name=U, below}, bend left=40, "a'"]
      & X \\
    & P \arrow[u, "p"']
    \arrow[from=U, to=1-2, Rightarrow, "\alpha", shorten <= 2pt, shorten >= 2pt]
  \end{tikzcd}
  \caption{The mediating triangle for the universal property of the pullback:
  $u\colon W \to P$ is the unique map with $p \circ u = a$. The double arrow
  $\alpha$ is a two-cell witnessing the comparison $a \Rightarrow a'$ between the
  two parallel maps $W \rightrightarrows X$.}
  \label{fig:universal}
\end{figure}

Taken together, \cref{fig:pullback,fig:universal} characterise the pullback up to
unique isomorphism: any two objects satisfying the cone condition and the
universal triangle are canonically isomorphic, since the mediating morphisms
compose to identities on both sides. This is the categorical content behind the
familiar set-level formula for $X \times_Z Y$.

\section{Conclusion}

The two examples, a short exact sequence in \cref{fig:ses} and a pullback square
in \cref{fig:pullback} with its universal triangle in \cref{fig:universal},
exhibit the principal arrow decorations of \texttt{tikz-cd}: labelled arrows,
hooked monomorphisms, two-headed epimorphisms, dashed mediating maps, the
pullback corner, bent arrows, and a two-cell. Each diagram is typeset directly
from its source, so the figures travel with the document without any external
image files.

\end{document}
